\section{符号说明}

全文共用的主要符号建议集中列于表\ref{tab:notation-demo}；仅在单问内部使用的局部参数可在首次出现处定义，避免符号表过长。

\begin{table}[H]
    \centering
    \caption{主要符号及单位}
    \label{tab:notation-demo}
    \begingroup
    \setlength{\extrarowheight}{2pt}
    \rowcolors{2}{gray!10}{white}
    \begin{tabularx}{0.97\textwidth}{
        >{\centering\arraybackslash}p{2.8cm}
        >{\centering\arraybackslash}X
        >{\centering\arraybackslash}p{2.5cm}}
        \toprule
        \rowcolor{CumcmLightBlue}
        \textbf{符号} & \textbf{含义} & \textbf{单位} \\
        \midrule
        $x$ & 厚度方向坐标 & m \\
        $t$ & 时间 & s \\
        $T_i(x,t)$ & 第$i$层衣料的温度分布 & $^\circ$C \\
        $\rho_i$ & 第$i$层衣料的密度 & kg/m$^3$ \\
        $c_i$ & 第$i$层衣料的比热容 & J/(kg$\cdot$K) \\
        $k_i$ & 第$i$层衣料的导热系数 & W/(m$\cdot$K) \\
        $d_i$ & 第$i$层衣料的厚度 & m \\
        $L$ & 三层衣料总厚度 & m \\
        $T_b$ & 人体恒定温度 & $^\circ$C \\
        $T_\infty$ & 环境温度 & $^\circ$C \\
        $T_{\mathrm{in}}(t)$ & 贴身侧温度（$T_1(0,t)$） & $^\circ$C \\
        $p(T)$ & 基线校正后的DSC放热功率 & W/kg \\
        $A_{\mathrm{DSC}}$ & DSC校正峰面积 & W$\cdot$K/kg \\
        $\xi_{\mathrm{DSC}}(T)$ & 由DSC数据定义的相对固化进度 & --- \\
        $L_{\mathrm{pcm}}$ & 功能层单位质量相变潜热 & kJ/kg \\
        $c_{2,\mathrm{eff}}(T)$ & 功能层有效比热容 & J/(kg$\cdot$K) \\
        $\beta_{\mathrm{DSC}}$ & DSC扫描速率假设 & K/min \\
        $h_b$ & 人体侧对流换热系数 & W/(m$^2\cdot$K) \\
        $h_e$ & 服装外表面对流换热系数 & W/(m$^2\cdot$K) \\
        $t_{15},t_{10}$ & 贴身侧温度首次降至$15^\circ$C、$10^\circ$C的时刻 & s \\
        $n$ & 新增最外层涂层数 & 层 \\
        $t_W(n)$ & 负重约束决定的站立时间 & s \\
        $t_{\mathrm{eff}}$ & 有效持续时间（热安全与负重约束的较小者） & s \\
        $\alpha$ & 功能层放热能力倍率 & --- \\
        $r_{\min}$ & 最小放热能力提高比例 & \% \\
        \bottomrule
    \end{tabularx}
    \endgroup
\end{table}
\FloatBarrier